\section{Setup: bias, transport, and the two dialects}
\label{sec:setup}

This section fixes the objects used throughout, proves the elementary
\emph{transport identity}, records the four equivalent forms of a
hypothetical Frankl counterexample, and introduces the two technical
devices the program runs on: the intersection-closed dialect and the
Walsh (harmonic) decomposition.

\subsection{Union-closed families and the bias}
\label{ssec:bias}

Fix a finite ground set $U$ with $|U|=n$; we identify $U$ with
$[n]=\{1,\dots,n\}$ and write $\powset$ for its power set, regarded as
the vertex set of the Boolean cube $\cube$.

\begin{definition}
\label{def:uc-family}
A family $\F\subseteq\powset$ is \emph{union-closed} if $A,B\in\F$
implies $A\cup B\in\F$. We exclude the degenerate cases
$\F=\emptyset$ and $\F=\{\emptyset\}$, and (harmlessly, by passing to
$\bigcup\F$) assume $U=\bigcup\F\in\F$. For $x\in U$ write
$\F_x=\{A\in\F:x\in A\}$ and $\F_{\bar x}=\F\setminus\F_x$. The
\emph{abundance} of $x$ is $a(x)=|\F_x|/|\F|$ and the \emph{bias} is
\[
  \bias_x \;=\; \bias_x(\F)\;=\;|\F_x|-|\F_{\bar x}|\;\in\;\Z.
\]
\end{definition}

In bias language Conjecture~\ref{conj:frankl} reads: every union-closed
$\F$ has a coordinate with $\bias_x\ge0$. We call such a coordinate
\emph{abundant}; a \emph{Frankl counterexample} is a union-closed
family with $\bias_x<0$ for every $x$.

\subsection{The transport identity}
\label{ssec:transport}

The $n$ coordinate involutions $\sigma_x\colon A\mapsto A\bigtriangleup
\{x\}$ generate the cube's structure group $(\Z/2)^n$ acting on
$\cube$. Restricted to a union-closed $\F$, each $\sigma_x$ is only a
\emph{partial} matching --- it pairs $A$ with $A\setminus\{x\}$ exactly
when both lie in $\F$. This partiality is the program's central object.

\begin{definition}
\label{def:transport}
The \emph{$x$-transport matching} is
$M_x=\{\{A,A\setminus\{x\}\}:A\in\F_x,\ A\setminus\{x\}\in\F\}$, the
direction-$x$ edges of the induced subgraph $G(\F)\subseteq\cube$. The
\emph{stuck sets} are
\[
  \F^{\stuck}_x=\{A\in\F_x:A\setminus\{x\}\notin\F\},\qquad
  \F^{\stuck}_{\bar x}=\{A'\in\F_{\bar x}:A'\cup\{x\}\notin\F\};
\]
these are the members \emph{needing} $x$, respectively \emph{rejecting}
$x$. The matched parts are $\F^{\matched}_x=\F_x\setminus\F^{\stuck}_x$
and $\F^{\matched}_{\bar x}=\F_{\bar x}\setminus\F^{\stuck}_{\bar x}$.
\end{definition}

\begin{proposition}[Transport identity]
\label{prop:transport-identity}
For every union-closed $\F$ and every $x\in U$,
\[
  \bias_x \;=\; |\F^{\stuck}_x|-|\F^{\stuck}_{\bar x}|.
\]
\end{proposition}

\begin{proof}
The matching $M_x$ pairs $\F^{\matched}_x$ bijectively with
$\F^{\matched}_{\bar x}$, so $|\F^{\matched}_x|=|\F^{\matched}_{\bar
x}|=|M_x|$. Hence $|\F_x|=|M_x|+|\F^{\stuck}_x|$ and
$|\F_{\bar x}|=|M_x|+|\F^{\stuck}_{\bar x}|$; subtract.
\end{proof}

Proposition~\ref{prop:transport-identity} says the bias is carried
\emph{entirely} by the stuck sets: the part of $\F$ on which the
coordinate toggle is defined contributes nothing. This is the precise
sense in which Frankl is about the \emph{failure} of the local symmetry
$\sigma_x$ to extend, and it motivates the whole program.

\begin{remark}
\label{rem:asymmetry}
The two sides of a coordinate are structurally \emph{different}
objects: $\F_x$ is a filter and is itself union-closed, whereas
$\F_{\bar x}$ is an order ideal and is a union-closed family on
$U\setminus\{x\}$. This asymmetry --- absent in the symmetric
$1/3$--$2/3$ setting --- is what the program exploits: union closure
\emph{orients} the obstruction. We return to it in
Section~\ref{ssec:uc1-uc2}.
\end{remark}

\subsection{Four equivalent forms of the obstruction}
\label{ssec:obstruction-forms}

A Frankl counterexample admits four equivalent descriptions, each the
entry point for a different technique. Let $\widehat{\mathbf
1_{\F}}(\{x\})$ denote the level-$1$ Walsh--Fourier coefficient of the
indicator $\mathbf1_\F\colon\powset\to\{0,1\}$ (see
Section~\ref{ssec:walsh}), and let $c(\F)\in[0,1]^n$ be the centroid,
$c(\F)_x=a(x)$.

\begin{proposition}
\label{prop:four-forms}
For a union-closed family $\F$ the following are equivalent.
\begin{enumerate}[label=\textup{(O\arabic*)},leftmargin=3em]
\item \textup{(combinatorial)} $a(x)<\tfrac12$ for every $x$;
\item \textup{(transport)}
  $|\F^{\stuck}_x|<|\F^{\stuck}_{\bar x}|$ for every $x$;
\item \textup{(harmonic)}
  $\widehat{\mathbf1_{\F}}(\{x\})>0$ for every $x$, using
  $\bias_x=-2^n\,\widehat{\mathbf1_{\F}}(\{x\})$;
\item \textup{(centroidal)} $c(\F)\in[0,\tfrac12)^n$.
\end{enumerate}
\end{proposition}

\begin{proof}
(O1)$\Leftrightarrow$(O2) is Proposition~\ref{prop:transport-identity};
(O1)$\Leftrightarrow$(O3) is the standard expansion
$\widehat{\mathbf1_\F}(\{x\})=2^{-n}\sum_A\mathbf1_\F(A)\,(-1)^{[x\in
A]}=2^{-n}(|\F_{\bar x}|-|\F_x|)$; (O1)$\Leftrightarrow$(O4) is the
definition of $a(x)$.
\end{proof}

Form (O4) is the slogan: \emph{Frankl says the centroid of a
union-closed family cannot lie strictly inside the small corner of the
cube.} Form (O3) exposes a feature with no $1/3$--$2/3$ analogue --- the
bias \emph{is}, up to scale and sign, a level-$1$ Fourier coefficient.
The harmonic form drives the entropy method~\cite{gilmer,sawin} and,
in a refined way, the entire cohomological arc.

\subsection{The intersection-closed dialect}
\label{ssec:flip}

The cohomological arc works with the order-dual presentation.

\begin{definition}
\label{def:int-closed}
A family $\F\subseteq\powset$ is \emph{intersection-closed} if
$A,B\in\F$ implies $A\cap B\in\F$. The \emph{flip}
$\tau(A)=[n]\setminus A$ sends a family $\F$ to $\F^c=\tau(\F)$.
\end{definition}

\begin{lemma}
\label{lem:flip}
$\F$ is union-closed if and only if $\F^c$ is intersection-closed;
moreover $|\F^c|=|\F|$ and $\bias_x(\F^c)=-\bias_x(\F)$.
\end{lemma}

\begin{proof}
$\tau$ is an order-reversing involution carrying $\cup$ to $\cap$;
$\tau(A)\ni x\iff x\notin A$, so $\bias_x$ changes sign.
\end{proof}

Hence Conjecture~\ref{conj:frankl} is equivalent to its
intersection-closed form: \emph{every intersection-closed family
$\F\neq\emptyset,\{[n]\}$ has a coordinate $x$ with $\bias_x\le0$}
(a \emph{rare} coordinate). The cohomological arc
(Sections~\ref{sec:cohomology}--\ref{sec:obstruction}) uses this
dialect because intersection-closed families form a category under
\emph{trace} maps (Definition~\ref{def:cn-category}) and because the
relevant symmetry group acts more transparently on it. The
combinatorial arc (Section~\ref{sec:combinatorial}) stays in the
union-closed dialect. Throughout, ``Frankl'' refers to either form via
Lemma~\ref{lem:flip}; we name the dialect when it matters.

\subsection{The Walsh decomposition}
\label{ssec:walsh}

The cube carries the regular representation of $(\Z/2)^n$, and its
irreducible characters are the Walsh characters.

\begin{definition}
\label{def:walsh}
For $S\subseteq[n]$ the \emph{Walsh character} $\walsh{S}\colon
\powset\to\{\pm1\}$ is $\walsh{S}(A)=(-1)^{|S\cap A|}$. The
characters $\{\walsh S\}_{S\subseteq[n]}$ are the $2^n$ irreducible
characters of $(\Z/2)^n$; under the coordinate toggle $\sigma_y$ one
has $\sigma_y\cdot\walsh S=(-1)^{[y\in S]}\walsh S$, and the
relabelling action of $\Sym_n$ permutes them by
$\pi\cdot\walsh S=\walsh{\pi(S)}$. We call $|S|$ the \emph{Walsh level}
of $\walsh S$; $\walsh{[n]}$ (every toggle acts by $-1$) is the
\emph{top} character.
\end{definition}

Two facts about the Walsh decomposition are used repeatedly and are
worth isolating now, because the audit in Section~\ref{sec:proof-state}
turns on them.

\begin{lemma}[Equivariant maps preserve isotype]
\label{lem:equivariant-preserves}
Let $V,W$ be $\Q$-representations of the finite abelian group
$(\Z/2)^n$ and let $f\colon V\to W$ be a $(\Z/2)^n$-equivariant
$\Q$-linear map. Then $f$ sends the $\walsh S$-isotypic subspace of $V$
into the $\walsh S$-isotypic subspace of $W$, for every $S$.
\end{lemma}

\begin{proof}
If $\sigma\cdot v=\walsh S(\sigma)\,v$ for all $\sigma$, then
$\sigma\cdot f(v)=f(\sigma\cdot v)=\walsh S(\sigma)\,f(v)$. (This is
Maschke / Schur for a finite abelian group: every representation
splits canonically into character isotypes, and an equivariant map is
block-diagonal for that splitting.)
\end{proof}

\begin{lemma}[Only multiplication shifts Walsh level]
\label{lem:cup-shifts}
Pointwise multiplication by a Walsh character, $v\mapsto\walsh U\cdot
v$, carries the $\walsh S$-isotype to the $\walsh{S\bigtriangleup
U}$-isotype, and is the \emph{only} natural operation that changes
Walsh level: it is not $(\Z/2)^n$-equivariant (it is $\walsh
U$-twisted-equivariant), so by Lemma~\ref{lem:equivariant-preserves}
no additive equivariant map can shift the level.
\end{lemma}

\begin{proof}
$\walsh S\cdot\walsh U=\walsh{S\bigtriangleup U}$, and
$\sigma_y\cdot(\walsh U\cdot v)=(-1)^{[y\in U]}\walsh U\cdot(\sigma_y
\cdot v)$.
\end{proof}

Lemmas~\ref{lem:equivariant-preserves} and~\ref{lem:cup-shifts} are
elementary, but they are exactly the mechanism on which the
obstruction-class landing ($Y_2$, Section~\ref{ssec:y2}) is sound and
the abandoned competing claim is impossible. We have stated them here
so that the audit can cite them without re-deriving them.
